\documentclass[11pt,a4paper]{article}

\usepackage[T1]{fontenc}
\usepackage[utf8]{inputenc}
\usepackage{lmodern}
\usepackage{microtype}
\usepackage[margin=24mm]{geometry}
\usepackage{amsmath,amssymb,mathtools}
\usepackage{booktabs,tabularx}
\usepackage{enumitem}
\usepackage{xcolor}
\usepackage{listings}
\usepackage{hyperref}
\usepackage{fancyhdr}

\definecolor{nismoblue}{HTML}{164E63}
\definecolor{nismogreen}{HTML}{166534}
\definecolor{nismored}{HTML}{991B1B}
\definecolor{nismogray}{HTML}{F3F4F6}
\definecolor{codegray}{HTML}{4B5563}

\hypersetup{
  colorlinks=true,
  linkcolor=nismoblue,
  urlcolor=nismoblue,
  citecolor=nismoblue,
  pdftitle={The Statistically Specified s-rwalk Implementation in NISMO},
  pdfauthor={NISMO contributors}
}

\lstdefinestyle{nismopython}{
  language=Python,
  basicstyle=\ttfamily\small,
  keywordstyle=\color{nismoblue}\bfseries,
  commentstyle=\color{nismogreen},
  stringstyle=\color{nismored},
  numbers=left,
  numberstyle=\scriptsize\color{codegray},
  numbersep=8pt,
  backgroundcolor=\color{nismogray},
  frame=single,
  rulecolor=\color{gray!35},
  breaklines=true,
  showstringspaces=false,
  columns=fullflexible,
  keepspaces=true
}
\lstset{style=nismopython}

\newcommand{\NISMO}{\textsc{NISMO}}
\newcommand{\code}[1]{\texttt{\detokenize{#1}}}
\newcommand{\R}{\mathbb{R}}

\setlist{nosep,leftmargin=*}
\setlength{\parindent}{0pt}
\setlength{\parskip}{0.55em}
\renewcommand{\arraystretch}{1.18}

\pagestyle{fancy}
\fancyhf{}
\lhead{\NISMO{} statistically specified \code{s-rwalk}}
\rhead{\thepage}
\setlength{\headheight}{14pt}

\title{\textbf{The Statistically Specified \code{s-rwalk}\\
Implementation in \NISMO}}
\author{Implementation specification for \code{nismo} version
\code{0.1.0.dev3}}
\date{\today}

\begin{document}

\maketitle

\begin{abstract}
This document specifies the separate
\code{proposal_scheme="s-rwalk"} kernel implemented by \NISMO.  It follows the
actual Python execution path, including immutable configuration, eligible
survivor selection, construction of a regularized survivor covariance,
Gaussian random-walk proposals, the fixed-\(q_0\) Metropolis correction,
randomized plateau handling, adaptive scale updates, resource limits, and
recorded diagnostics.  The existing Dynesty-style \code{"rwalk"} and ensemble
\code{"en-rwalk"} implementations are outside the scope of this document and
are not reinterpreted here.
\end{abstract}

\tableofcontents
\newpage

\section{Scope and implementation locations}

This document concerns only \code{"s-rwalk"}.  It is a scalar-chain,
Gaussian-covariance Metropolis--Hastings replacement kernel.  It coexists with:

\begin{itemize}
  \item \code{"rwalk"}, the existing Dynesty-style ellipsoidal-ball walk; and
  \item \code{"en-rwalk"}, the split differential-evolution ensemble walk.
\end{itemize}

The implementation is divided among four modules:

\begin{itemize}
  \item \code{src/nismo/config.py}: public \code{SRWalkSettings} and the
        \code{"s-rwalk"} proposal-scheme literal;
  \item \code{src/nismo/mcmc.py}: \code{SRWalkSampler}, regularized covariance
        factor, Gaussian proposal evolution, fixed-\(q_0\) acceptance, and
        scale tuning;
  \item \code{src/nismo/sampler.py}: controller lifetime, dispatch, resource
        integration, live-point replacement, and per-iteration histories;
  \item \code{src/nismo/__init__.py}: public export of
        \code{SRWalkSettings}.
\end{itemize}

No external MCMC package is required.  The code uses the sampler's supplied
NumPy \code{Generator} exclusively.

\section{The constrained density targeted by the walk}

Let \(L(\theta)\) be the likelihood, \(\pi(\theta)\) the normalized prior,
and \(q_0(\theta)\) the fixed normalized importance density supplied as
\code{importance_morph}.  \NISMO{} defines

\begin{equation}
  \log \Psi_0(\theta)
  = \log L(\theta)+\log\pi(\theta)-\log q_0(\theta).
  \label{eq:psi}
\end{equation}

At one nested iteration, the discarded live point defines the threshold
\(\lambda\).  The replacement kernel must preserve

\begin{equation}
  p_\lambda(\theta)
  \propto q_0(\theta)
  \mathbf{1}\!\left[\log\Psi_0(\theta)>\lambda\right].
  \label{eq:target}
\end{equation}

Equation~\eqref{eq:target} is not the ordinary nested-sampling constrained
prior unless \(q_0=\pi\).  The likelihood and prior determine whether a point
passes the constraint through Equation~\eqref{eq:psi}; the Metropolis density
ratio itself is a ratio of the fixed importance density \(q_0\).

\section{Public configuration and exact resolution}

The user passes an immutable \code{SRWalkSettings} object:

\begin{lstlisting}
SRWalkSettings(
    n_steps=25,
    scale=None,
    facc=0.5,
    covariance_shrinkage=0.1,
    covariance_jitter=1.0e-10,
)
\end{lstlisting}

For model dimension \(D\), the run-level \code{SRWalkSampler} resolves the
values as follows.  Write \(S\) for the configured number of steps,
\(f_0\) for the resolved target acceptance fraction, \(\rho\) for covariance
shrinkage, and \(\epsilon\) for covariance jitter.

\begin{center}
\begin{tabularx}{\textwidth}{@{}l l X@{}}
\toprule
Field & Default or resolved value & Validation and meaning \\
\midrule
\code{n_steps}
& \(S=25\)
& Must be a positive integer.  Every successful replacement attempts exactly
  \(S\) scalar transitions. \\
\code{scale}
& \(s_0=2.38/\sqrt{D}\) when omitted
& An explicit value must be positive and finite.  It replaces the initial
  value but remains subject to subsequent adaptation. \\
\code{facc}
& \(f_0=\min(1,\max(1/S,f_{\rm input}))\)
& The input must be finite.  The controller clamps it to the same range used
  by the existing \code{rwalk} controller. \\
\code{covariance_shrinkage}
& \(\rho=0.1\)
& Must be finite and lie in \([0,1]\). \\
\code{covariance_jitter}
& \(\epsilon=10^{-10}\)
& Must be positive and finite.  It is both an additive diagonal
  regularizer and the deterministic eigenvalue floor. \\
\bottomrule
\end{tabularx}
\end{center}

The factor \(2.38/\sqrt{D}\) is a mixing heuristic, not part of the
correctness of the target density.  Each call to \code{NISMOSampler.run()}
constructs a fresh \code{SRWalkSampler}; scale adaptation does not leak from
one run call to the next.

\section{Threshold, starting state, and constraint ordering}

\subsection{The discarded point is not a chain state}

The current worst point is identified from the stored
\code{live_log_psi0} ordering, with the stored tie breaker included when the
randomized plateau policy is active.  That point supplies \(\lambda\), but
under strict ordering it satisfies equality and lies outside
Equation~\eqref{eq:target}.  It is therefore explicitly excluded from the
eligible start indices.

The implementation evaluates the shared constraint helper on the stored live
arrays, removes the worst index, and selects one eligible survivor uniformly:

\begin{equation}
  I_0\sim\operatorname{Uniform}(\mathcal{I}_{\rm eligible}).
\end{equation}

The complete cached state is copied from row \(I_0\):

\begin{equation}
  \left(\theta,\log L,\log\pi,\log q_0,\log\Psi_0,t\right).
\end{equation}

The starting state is not reevaluated.  If no eligible survivor exists, the
replacement fails with \code{insufficient_eligible_survivors}.

\subsection{Strict and randomized plateau policies}

Under \code{tie_policy="strict"}, a proposal passes when

\begin{equation}
  \log\Psi_0(\theta')>\lambda.
\end{equation}

Under \code{tie_policy="randomized_plateau"}, every state is augmented by
\(t\sim\operatorname{Uniform}(0,1)\).  The discarded state defines
\((\lambda,t_\lambda)\), and every parameter proposal receives a fresh
\(t'\).  Passage is lexicographic:

\begin{equation}
  \log\Psi_0(\theta')>\lambda
  \quad\text{or}\quad
  \left[\log\Psi_0(\theta')=\lambda\ \text{and}\ t'>t_\lambda\right].
  \label{eq:plateau}
\end{equation}

On acceptance, both \(\theta'\) and \(t'\) become current.  On either a
constraint rejection or an MH rejection, both current values are retained.

\section{Frozen survivor covariance geometry}

\subsection{Input survivor set}

Let the complete live matrix be

\begin{equation}
  \Theta=(\theta_1,\ldots,\theta_N)^\mathsf{T}\in\R^{N\times D}.
\end{equation}

For \code{s-rwalk}, the covariance input excludes the discarded worst row:

\begin{equation}
  X=\Theta_{-\mathrm{worst}}\in\R^{M\times D},
  \qquad M=N-1.
\end{equation}

This differs from standard \code{rwalk}, whose bounding ellipsoid is built
from the complete live set.  The \code{s-rwalk} covariance is calculated once
per attempted replacement and is not cached between replacements.  All
entries of \(X\) must be finite.

\subsection{Empirical covariance and shrinkage}

Let

\begin{equation}
  \bar{x}=\frac{1}{M}\sum_{i=1}^{M}x_i,
  \qquad
  Y_i=x_i-\bar{x}.
\end{equation}

The implementation computes

\begin{equation}
  C=\frac{Y^\mathsf{T}Y}{\max(M-1,1)}.
  \label{eq:empirical-covariance}
\end{equation}

The denominator in Equation~\eqref{eq:empirical-covariance} keeps the helper
defined even when only one survivor is available.  It also handles \(D=1\)
without a scalar covariance special case.

The regularized covariance is

\begin{equation}
  C_{\rm reg}
  =(1-\rho)C+\rho\,\operatorname{diag}(C)+\epsilon I_D.
  \label{eq:regularized-covariance}
\end{equation}

Shrinkage damps empirical cross-covariances while retaining marginal
variances.  The diagonal jitter ensures a positive baseline in directions
that the finite live set does not resolve.

\subsection{Deterministic positive-definite factor}

After explicit symmetrization, the implementation computes

\begin{equation}
  C_{\rm reg}=V\operatorname{diag}(\ell_1,\ldots,\ell_D)V^\mathsf{T}.
\end{equation}

Every eigenvalue is floored deterministically:

\begin{equation}
  \tilde{\ell}_j=\max(\ell_j,\epsilon).
\end{equation}

The returned proposal factor is

\begin{equation}
  L=V\operatorname{diag}(\sqrt{\tilde{\ell}_1},\ldots,
                          \sqrt{\tilde{\ell}_D}),
  \qquad
  LL^\mathsf{T}=V\operatorname{diag}(\tilde{\ell})V^\mathsf{T}.
  \label{eq:factor}
\end{equation}

This construction supports one dimension, fewer survivors than dimensions,
and rank-deficient survivor clouds.  A \code{NumericalInvariantError} is raised
only if the input, regularized matrix, eigendecomposition, or resulting factor
cannot be made finite and usable.

\subsection{Geometry lifetime}

The factor \(L\) is frozen for all \(S\) transitions in one replacement.  It
is not recomputed after accepted or rejected proposals.  The scale \(s_k\)
used by that replacement is frozen at the same time.  Only after a complete
walk is returned successfully is the controller scale adapted for the next
replacement.

\section{One Gaussian proposal}

At each transition, draw

\begin{equation}
  z\sim\mathcal{N}(0,I_D)
\end{equation}

and propose

\begin{equation}
  \theta'=\theta+s_kLz.
  \label{eq:proposal}
\end{equation}

Conditional on frozen \(L\) and \(s_k\), the increment is zero-mean Gaussian
with covariance

\begin{equation}
  \operatorname{Cov}(\theta'-\theta)
  =s_k^2LL^\mathsf{T}.
\end{equation}

Thus the parameter-space proposal is symmetric:

\begin{equation}
  Q(\theta'\mid\theta;L,s_k)
  =Q(\theta\mid\theta';L,s_k).
  \label{eq:symmetry}
\end{equation}

\NISMO{} proposes in the full physical parameter space.  It does not clip,
reflect, wrap, or repeatedly redraw a point at prior boundaries.  A point with
zero prior or likelihood simply fails the active pseudo-likelihood constraint.

\section{Evaluation and fixed-\texorpdfstring{\(q_0\)}{q0} Metropolis decision}

Every proposal is evaluated exactly once by \code{BatchEvaluator}, producing

\begin{equation}
  \left(\theta',\log L',\log\pi',\log q_0',\log\Psi_0'\right),
\end{equation}

where \(\log q_0'\) always comes from the original fixed
\code{importance_morph}.  The adaptive scale and survivor covariance never
replace this density.

First, the implementation draws the proposed tie value and tests the active
constraint.  A constraint failure is rejected without an MH uniform draw.  If
the proposal passes, symmetry in Equation~\eqref{eq:symmetry} gives

\begin{equation}
  \log\alpha
  =\min\!\left(0,
    \log q_0(\theta')-\log q_0(\theta)
  \right).
  \label{eq:mh}
\end{equation}

The proposal is accepted when

\begin{equation}
  \log U<\log\alpha,
  \qquad U\sim\operatorname{Uniform}(0,1).
\end{equation}

Neither \(\log L'-\log L\), \(\log\pi'-\log\pi\), nor
\(\log\Psi_0'-\log\Psi_0\) is substituted for the fixed-\(q_0\) ratio.  An
accepted proposal replaces every cached current field.  A rejected proposal
changes none of them.

\section{Exact complete-walk algorithm}

For one nested replacement, the implemented sequence is:

\begin{enumerate}
  \item Validate all live-array shapes.
  \item Preflight resource limits for all \(S\) proposals.
  \item Build the eligible survivor indices and fail if the set is empty.
  \item Delete the worst row from the frozen live matrix and construct \(L\)
        from Equations~\eqref{eq:regularized-covariance}--\eqref{eq:factor}.
  \item Select one eligible survivor uniformly and copy its complete cached
        state without reevaluation.
  \item Freeze the controller's current scale \(s_k\).
  \item Repeat exactly \(S\) times:
    \begin{enumerate}
      \item check the wall-time deadline;
      \item draw \(z\), form Equation~\eqref{eq:proposal}, and evaluate it once;
      \item draw a fresh proposed tie value;
      \item test the active constraint;
      \item if valid, apply Equation~\eqref{eq:mh};
      \item on acceptance, replace the complete current state.
    \end{enumerate}
  \item Return the actual final chain state, even when no proposal was
        accepted.
  \item Tune the scale using the complete walk's acceptance count.
\end{enumerate}

The following pseudocode emphasizes the execution ordering:

\begin{lstlisting}
eligible = eligible_survivor_indices(...)
survivors = delete(live_theta, worst)
factor = covariance_factor(
    survivors,
    shrinkage=sampler.covariance_shrinkage,
    jitter=sampler.covariance_jitter,
)
current = copy(uniform_choice(eligible))
scale_used = sampler.scale

for _ in range(sampler.n_steps):
    z = rng.standard_normal(ndim)
    proposed = evaluate(
        current.theta + scale_used * factor @ z
    )
    proposed.tie = rng.random()

    if not passes_constraint(proposed):
        continue
    n_valid += 1

    log_alpha = min(
        0.0,
        proposed.log_q0 - current.log_q0,
    )
    if log(rng.random()) < log_alpha:
        current = proposed
        n_accepted += 1

sampler.record_completed_walk(
    accept=n_accepted,
    scale=scale_used,
)
return current
\end{lstlisting}

The matrix multiplication in the implementation is explicitly
\code{factor @ z}; the proposal is therefore distributed according to the
regularized covariance in Equation~\eqref{eq:factor}.

\section{Adaptive scale between replacements}

Let \(A_k\) and \(R_k\) be the accepted and rejected transitions recorded for
one completed \code{s-rwalk} replacement.  In the serial implementation,

\begin{equation}
  A_k+R_k=S,
  \qquad
  f_k=\frac{A_k}{A_k+R_k}.
\end{equation}

Here ``accepted'' means acceptance after both the active constraint and the
fixed-\(q_0\) MH test.  Constraint failures and valid-but-MH-rejected points
both contribute to \(R_k\).  The scale update is the same acceptance-target
recursion used by the existing \code{rwalk} controller:

\begin{equation}
  s_{k+1}
  =s_k\exp\!\left(\frac{f_k-f_0}{D f_0}\right).
  \label{eq:tune}
\end{equation}

Acceptance above \(f_0\) expands the next replacement's proposal scale;
acceptance below \(f_0\) shrinks it.  The scale used inside the completed walk
is never changed retroactively.  After Equation~\eqref{eq:tune}, the
controller's acceptance and rejection counters reset to zero.

Because an explicit \code{scale} is an initial value, not a fixed-scale mode,
Equation~\eqref{eq:tune} also applies when the user supplies \code{scale}.

\section{Resource limits and failure semantics}

Before choosing a start or constructing covariance geometry, \NISMO{} rejects
the complete replacement without new likelihood calls when:

\begin{itemize}
  \item \(S>\code{max_proposals_per_replacement}\);
  \item the remaining run-wide likelihood budget cannot accommodate all
        \(S\) calls; or
  \item the wall-time deadline has already elapsed.
\end{itemize}

The wall-time deadline is checked again before every scalar transition.  If it
expires inside a walk, the attempt returns no draw, the worst live point is
left unchanged, and the shortened chain is not used.  The controller is tuned
only after a successful complete walk.  A successful replacement has

\begin{equation}
  n_{\rm proposed}
  =n_{\rm likelihood\ calls}
  =n_{\rm completed}
  =S,
\end{equation}

where likelihood calls in this equality refer only to calls made by the
replacement.  Initial live-set evaluation is additional run-wide work.

\section{Recorded diagnostics}

For each completed \code{s-rwalk} replacement, \code{RunHistory} records:

\begin{center}
\begin{tabularx}{\textwidth}{@{}l X@{}}
\toprule
Field & \code{s-rwalk} meaning \\
\midrule
\code{constraint_pass_fraction}
& Constraint passes divided by all \(S\) generated candidates. \\
\code{mh_acceptance_fraction}
& Final MH acceptances divided by all \(S\) generated candidates, not only
  by constraint-passing candidates. \\
\code{mcmc_accepted}
& Count of accepted transitions in the completed chain. \\
\code{mcmc_moved}
& One if at least one transition was accepted, otherwise zero. \\
\code{mcmc_completed}
& Exactly \(S\). \\
\code{acceptance_fraction}
& Cumulative nested replacements divided by all proposals.  This retains the
  sampler-wide efficiency meaning and is not the MH rate. \\
\bottomrule
\end{tabularx}
\end{center}

The current adapted scale and covariance factor are controller-local state and
are not added as public \code{RunHistory} arrays.

\section{Why the kernel preserves the desired density}

Conditional on the frozen survivor geometry \(L\) and scale \(s_k\), the
Gaussian proposal is symmetric by Equation~\eqref{eq:symmetry}.  For two
states inside the permitted constraint, the target ratio from
Equation~\eqref{eq:target} is

\begin{equation}
  \frac{p_\lambda(\theta')}{p_\lambda(\theta)}
  =\frac{q_0(\theta')}{q_0(\theta)},
\end{equation}

which yields Equation~\eqref{eq:mh}.  A proposal outside the active constraint
has target density zero and is rejected.  The same argument applies to the
augmented \((\theta,t)\) state under Equation~\eqref{eq:plateau}, because the
tie variable has a uniform density that cancels in the MH ratio.

The covariance and scale remain fixed for the entire chain.  Covariance is
recomputed and scale is adapted only between completed replacements, so the
proposal mechanics do not change in response to accepted or rejected states
within the active walk.

The eligible starting state is copied from a surviving live point already in
the constrained target.  Consequently, any finite number of valid
Metropolis transitions preserves the constrained distribution.  This is an
invariance statement, not an independence or mixing guarantee.

\section{Distinction from the existing \texorpdfstring{\code{rwalk}}{rwalk}}

The public schemes are intentionally separate:

\begin{center}
\begin{tabularx}{\textwidth}{@{}l X X@{}}
\toprule
Property & \code{rwalk} & \code{s-rwalk} \\
\midrule
Settings object
& \code{RWalkSettings}
& \code{SRWalkSettings} \\
Displacement
& Uniform in a transformed unit ball
& Multivariate Gaussian \\
Geometry
& Single ellipsoid bounding the complete live set
& Regularized covariance of survivors only \\
Geometry lifetime
& Cached for approximately \code{walks * n_live} calls
& Recomputed once per replacement \\
Default transitions
& \(D+20\)
& 25 \\
Initial scale
& 1
& \(2.38/\sqrt{D}\) \\
Scale adaptation
& Acceptance-target recursion
& The same acceptance-target recursion \\
MH target
& Constrained fixed \(q_0\)
& Constrained fixed \(q_0\) \\
\bottomrule
\end{tabularx}
\end{center}

Selecting \code{"s-rwalk"} does not route through
\code{draw_rwalk_constrained}, does not use the bounding-ellipsoid cache, and
does not modify the behavior of \code{"rwalk"}.

\section{Mixing limitations}

A finite-length Metropolis kernel can preserve the correct density while
remaining highly correlated with its selected starting survivor.  Poor
mixing may appear as low constraint-pass fractions, low MH acceptance, or
frequent \code{mcmc_moved=0} outcomes.  Users should:

\begin{itemize}
  \item compare evidence and posterior summaries across repeated complete
        runs and seeds;
  \item test sensitivity to \code{n_steps}, \code{facc}, covariance
        regularization, and \code{n_live};
  \item inspect both constraint-pass and unconditional MH acceptance
        histories; and
  \item compare against independent constrained rejection when practical.
\end{itemize}

\NISMO{} does not force an accepted transition, discard unchanged outputs,
select the highest-likelihood visited state, or choose a state conditional on
movement.  Each of those operations would alter the intended target.

\section{Usage}

Default geometry and scale initialization with a longer chain:

\begin{lstlisting}
from nismo import NISMOSampler, SRWalkSettings

sampler = NISMOSampler(
    model=model,
    importance_morph=importance_morph,
    proposal_scheme="s-rwalk",
    srwalk_settings=SRWalkSettings(
        n_steps=50,
        facc=0.5,
    ),
    n_live=200,
    rng=42,
)

result = sampler.run(
    dlogz=0.1,
    max_proposals_per_replacement=100,
)
\end{lstlisting}

Explicit initial scale and covariance regularization:

\begin{lstlisting}
sampler = NISMOSampler(
    model=model,
    importance_morph=importance_morph,
    proposal_scheme="s-rwalk",
    srwalk_settings=SRWalkSettings(
        n_steps=40,
        scale=0.75,
        facc=0.4,
        covariance_shrinkage=0.2,
        covariance_jitter=1.0e-8,
    ),
    n_live=200,
    rng=42,
)
\end{lstlisting}

The second example starts at scale 0.75, then adapts that value after each
completed replacement.  The hard proposal limit passed to
\code{NISMOSampler.run()} must be at least \code{n_steps}.

Useful inspection:

\begin{lstlisting}
print(result.history.constraint_pass_fraction)
print(result.history.mh_acceptance_fraction)
print(result.history.mcmc_accepted)
print(result.history.mcmc_moved)
print(result.history.mcmc_completed)
\end{lstlisting}

\section{Validation coverage}

The implementation is covered at three levels:

\begin{itemize}
  \item deterministic unit tests validate setting rejection, default and
        explicit scale initialization, the adaptation recursion, uniform
        eligible starts, survivor-only frozen covariance, exact proposal and
        likelihood counts, all-rejected chains, and preflight failures;
  \item integration tests validate public sampler dispatch, histories,
        resource behavior, and fixed-seed reproducibility; and
  \item statistical tests validate stationarity for a truncated normal target,
        a correlated Gaussian with a non-axis-aligned constraint, and repeated
        end-to-end evidence runs.
\end{itemize}

The statistical tests are specifically sensitive to the fixed-\(q_0\) MH
correction: accepting every proposal merely because it passes the constraint
does not preserve the tested truncated-normal distribution.

\section{References and runtime citation behavior}

The Gaussian random-walk Metropolis mechanism is the symmetric-proposal case
of the Metropolis--Hastings algorithm.  The \(2.38/\sqrt{D}\) initialization
is a conventional high-dimensional Gaussian-target scaling heuristic; it is
not required for detailed balance.  The \NISMO{} fixed-importance target and
nested replacement context follow the project's mathematical contract.

At runtime, \code{sampler.citations} for \code{"s-rwalk"} currently returns
the same Skilling (2006) nested-sampling method entry exposed by
\code{"rwalk"}.  The additional references below explain the general MH and
scaling context; they are not additional runtime citation entries.

\begin{thebibliography}{9}

\bibitem{metropolis1953}
N. Metropolis, A. W. Rosenbluth, M. N. Rosenbluth, A. H. Teller, and
E. Teller,
``Equation of state calculations by fast computing machines,''
\emph{The Journal of Chemical Physics}, vol.~21, no.~6, pp.~1087--1092, 1953.

\bibitem{hastings1970}
W. K. Hastings,
``Monte Carlo sampling methods using Markov chains and their applications,''
\emph{Biometrika}, vol.~57, no.~1, pp.~97--109, 1970.

\bibitem{roberts1997}
G. O. Roberts, A. Gelman, and W. R. Gilks,
``Weak convergence and optimal scaling of random walk Metropolis
algorithms,''
\emph{The Annals of Applied Probability}, vol.~7, no.~1, pp.~110--120, 1997.

\bibitem{skilling2006}
J. Skilling,
``Nested sampling for general Bayesian computation,''
\emph{Bayesian Analysis}, vol.~1, no.~4, pp.~833--859, 2006.
\url{https://projecteuclid.org/euclid.ba/1340370944}

\end{thebibliography}

\end{document}
